\pdfbookmark[1]{Auflösung quadratischer Gleichungen. 1840.}{quadratisch1840}\par\Needspace{10\baselineskip}\begin{center}\rule{0.7\linewidth}{0.4pt}\par \textsc{Vega} und \textsc{Hülsse.} Sammlung mathematischer Tafeln. 1840.\par\rule{0.7\linewidth}{0.4pt}\end{center}

\noindent\textit{Lösung quadratischer Gleichungen in der Form, dass nicht die Koeffizienten}
\textit{der Gleichung selbst, sondern ihre Logarithmen gegeben sind und man auch}
\textit{nicht die Wurzeln selbst (oder eine von ihnen), sondern ihre Logarithmen}
\textit{zur weiteren Verwendung benötigt.}

Die Durchführung dieser Aufgabe allein mithilfe der gewöhnlichen Logarithmentafeln erfordert notwendig viermaliges Nachschlagen. Zweimaliges Nachschlagen genügt, wenn man entweder die trigonometrischen Logarithmentafeln oder die Tafeln für Logarithmen von Summen und \SourceLeaf{265} Differenzen verwendet. Einmaliges Nachschlagen genügt, wenn die letztere Tafel wie hier um die Zusatzspalten $D$, $E$ und $F$ erweitert ist.

Die vierte Spalte, $D$, enthält $B+C$, die fünfte, $E$, enthält $A+C$,
die sechste, $F$, enthält $B-A$. Nehmen wir an, dass die Zahlen selbst, zu denen die Logarithmen $A$, $B$, $C$ gehören, jeweils $a$, $b$, $c$ sind. Dann enthielte also die vierte Spalte die Logarithmen von $bc$, die fünfte die von $ac$, die
sechste die von $\frac{b}{a}$.

Es sei noch darauf hingewiesen, dass die Logarithmen in der sechsten Spalte $F$, die bis zu $A=0{,}208$ positiv sind, von $A=0{,}209$
an negativ werden (S.\ 640). Es ist ziemlich gleichgültig, ob sie in dieser negativen Form oder wie hier durch ihre Komplemente angegeben werden. So könnte zum Beispiel für $A=0{,}367$ unter $F$ stehen $-0{,}21180$ oder, wie hier, $9{,}78820$.

Bei der Anwendung dieser Zusatzspalten auf die Lösung der quadratischen Gleichung
\[
pxx+qx+r=0
\]
selbst sind vier verschiedene Fälle zu unterscheiden:
\begin{enumerate}
\item[I.] $p$ und $r$ haben gleiche Vorzeichen und $\frac{qq}{pr}$ ist nicht kleiner als $4$;
\item[II.] $p$ und $r$ haben gleiche Vorzeichen und $\frac{qq}{pr}$ ist kleiner als $4$;
\item[III.] $p$ und $r$ haben entgegengesetzte Vorzeichen und $-\frac{pr}{qq}$ ist größer als $2$;
\item[IV.] $p$ und $r$ haben entgegengesetzte Vorzeichen und $-\frac{pr}{qq}$ ist kleiner als $2$.
\end{enumerate}
Der Fall, in dem $-\frac{pr}{qq}=2$ ist, kann sowohl zu III.\ als auch zu IV.\ gerechnet werden.

\Needspace{24\baselineskip}
Im Fall II.\ sind die Wurzeln imaginär; in den übrigen Fällen erhält man jede der beiden Wurzeln auf zwei Arten. Das folgende Schema macht alles leicht überschaubar. Dabei ist
\[
\frac{q}{p}=h,\quad \frac{r}{q}=g
\]
gesetzt, teils zur Abkürzung, teils weil die Rechnung tatsächlich in dieser Form am bequemsten ausgeführt wird.

\begin{center}
\begingroup\renewcommand{\arraystretch}{1.7}\begin{tabular}{rlll}
 & & \emph{Erste Wurzel.} & \emph{Zweite Wurzel.} \\
I. & $+\frac{h}{g}=bc=d$ & $-\frac{h}{b}=-gc$ & $-gb=-\frac{h}{c}$ \\
II. & $+\frac{h}{g}<4$ & Imaginär. & Imaginär. \\
III. & $-\frac{g}{h}=ac=e$ & $+ha=-\frac{g}{c}$ & $+\frac{g}{a}=-hc$ \\
IV. & $-\frac{g}{h}=\frac{b}{a}=f$ & $+\frac{h}{a}=-\frac{g}{b}$ & $+ga=-hb$
\end{tabular}\endgroup
\end{center}
\par
\SourceLeaf{266}
Die Beweise dieser Regeln wird jeder leicht selbst entwickeln können. Zum Überfluss sei noch ein Anwendungsbeispiel angegeben.

\begin{center}
\begin{tabular}{rl}
Es sei gegeben & $\log p=0{,}69897$ \\
 & $\log q=0{,}84510$ \\
 & $\log r=0{,}77815\,\mathrm{n}$ \\
Also & $\log h=0{,}14613$ \\
 & $\log g=9{,}93305\,\mathrm{n}$ \\
 & $\log(-\frac{g}{h})=9{,}78692$
\end{tabular}
\end{center}

Man erkennt unmittelbar, dass hier der Fall IV.\ vorliegt und daher $9{,}78692$ in der sechsten Spalte unter $B-A$ oder $F$ zu suchen ist. Von der Tafel (S.\ 643) wäre hier nur Folgendes nötig:
\begin{center}
\begin{tabular}{c|c|c}
$A$ & $B$ & $B-A$ oder $F$ \\
$0{,}367$ & $0{,}15520$ & $9{,}78820$ \\
$0{,}368$ & $0{,}15489$ & $9{,}78689$
\end{tabular}
\end{center}
Zu $F=9{,}78692$ gehört also
\begin{center}
\begin{tabular}{l}
$A=0{,}36798=\log a$, woraus $\log\frac{h}{a}=9{,}77815$, $\log ga=0{,}30103\,\mathrm{n}$ oder \\
$B=0{,}15490=\log b$, woraus $\log(-\frac{g}{b})=9{,}77815$, $\log(-hb)=0{,}30103\,\mathrm{n}$.
\end{tabular}
\end{center}

